\subsection{Incomplete polynomial oracles}
\label{sec:incomplete-oracles}

Recall at the end of a PIOR that the verifier outputs a new instance $\inp' \gets \verifier\bigl(\idx,\inp,\provermsg_1,\verifiermsg_1,\dots,\provermsg_k,\verifiermsg_k\bigr)$, and any oracles therein are specified implicitly in terms of oracles given in $\inp$ and $\provermsg_1,\dots,\provermsg_k$. For example, the new instance could contain an oracle $\oracle{f'} = \oracle{f} + \verifiermsg_1 \cdot \oracle{\provermsg_1}$, for some $\oracle{f}$ included in the old instance. Assuming $\verifiermsg_1$ and the coefficients of $f$ and $\provermsg_1$ all reside over the same ring, this new polynomial oracle is complete, in that it can be evaluated at any set of points in that ring.

In Zinc$+$ the verifier outputs a new instance with polynomial oracles of the form $\psi_\verifiermsg([[f(\psi_{\verifiermsg}^{-1}{(\cdot)})]])$, where we can always compute $v' = [[f(\psi_{\verifiermsg}^{-1}{(\cdot)})]]$, but for malicious provers $v'$ may lie outside the domain $\psi_\verifiermsg$ is defined on. We model this via the (possibly) \emph{incomplete polynomial oracle} $\incompleteoracle{f'} = \psi_\verifiermsg([[f(\psi_{\verifiermsg}^{-1}{(\cdot)})]])$, which either outputs a correct evaluation or $\perp$ if an undefined operation is encountered during the evaluation, and which we syntactically distinguish from complete polynomial oracles using braces rather than brackets.

We can modify any PIOP to work with possibly incomplete oracles as follows. If the oracle outputs $\perp$, the verifier automatically rejects, and otherwise the protocol proceeds normally. If we can show that during the execution of the modified PIOP incomplete oracles that are not complete output $\perp$ whp, then whenever the verifier accepts we can assume whp that the possibly incomplete oracle is in fact complete. For the possibly incomplete oracle $\incompleteoracle{f'}$ used in Zinc$+$, we show that when queried on a random point it outputs $\perp$ whp. This means that the very weak requirement that the PIOP queries $\incompleteoracle{f'}$ on a random point suffices for the security of our reduction, and in practice this does not rule out any public-coin PIOPs of interest.